\section{Linear symmetries of the level equations}

This section distinguishes preservation of the individual level equations from
preservation of the scheme they define.  The distinction is necessary before a
fixed-point argument can be applied to their Grassmannians.

For a matrix $t=(t_{ij})\in\operatorname{Mat}_n(\mathbb C)$, let
\[
 S_k(t)=\sum_{\substack{R\subseteq\{1,\ldots,n\}\\|R|=k}}
        \ \sum_{\substack{f:R\hookrightarrow\{1,\ldots,n\}}}
             \prod_{i\in R}t_{i,f(i)}.
\]
In particular $S_1$ is the entry sum and $S_n$ is the permanent.

\begin{proposition}\label{prop:conformal-layers}
Let $n\geq3$.  Suppose that $T\in\operatorname{GL}(\operatorname{Mat}_n(\mathbb C))$
and that for each $k$ there is $c_k\in\mathbb C^*$ such that
$S_k(Tt)=c_kS_k(t)$ for all $t$.  Then $T$ is a nonzero scalar times a row
permutation, a column permutation, and possibly transposition.  Conversely every
such transformation has the asserted property.  The scalars satisfy $c_k=c_1^k$.
\end{proposition}

\begin{proof}
At $k=n$ the hypothesis says that $T$ preserves the permanent up to scalar.
Marcus and May's main theorem \cite[Section~2, Theorem, p.~179,
(2.2)--(2.4)]{MarcusMay1962} therefore writes $T(t)$ as $DPtQL$ or $DPt^{\mathsf T}QL$,
where $P,Q$ are permutation matrices and $D,L$ are invertible diagonal matrices.
Undo the permutations and transposition.  The entry $t_{ij}$ is multiplied by
$d_i l_j$.  Comparing its coefficient in $S_1(Tt)=c_1S_1(t)$ gives
$d_i l_j=c_1$ for every pair $i,j$.  Thus all entry multipliers are the same.
The converse and the relation $c_k=c_1^k$ follow from homogeneity.
\end{proof}

The scalar subgroup fixes every vector subspace.  Thus this group supplies no
nontrivial torus action on a Grassmannian.  Preserving the ideal of all equations
is a weaker condition: a generator can change by multiples of lower-degree
generators.  The following finite calculation treats that weaker condition at
order four.

\begin{theorem}\label{thm:scheme-torus-four}
Let $H=\{t\in\operatorname{Mat}_4(\mathbb C):S_1(t)=0\}$ and let
\[
 I=(S_2|_H,S_3|_H,S_4|_H)\subseteq\mathbb C[H].
\]
The Lie algebra of the linear stabilizer of $I$ consists of scalar endomorphisms
of $H$.  Every algebraic torus acting linearly on $H$ and preserving $I$ therefore
acts by scalars, and fixes every point of every Grassmannian of subspaces of $H$.
\end{theorem}

\begin{proof}
Use the fifteen entries other than $t_{44}$ as coordinates
$x_1,\ldots,x_{15}$ on $H$, in row-major order, and put
$t_{44}=-\sum_{i=1}^{15}x_i$.  Write $Q=S_2|_H$ and $C=S_3|_H$.
The homogeneous pieces of $I$ in degrees two and three are
\[
 I_2=\mathbb C Q,\qquad I_3=\mathbb C C+QH^*.
\]
Consequently an infinitesimal endomorphism $A=(a_{ij})$ preserving $I$ must
satisfy
\begin{align}
 \sum_{i,j=1}^{15}a_{ij}x_j\frac{\partial Q}{\partial x_i}
     &=\lambda_2Q,\\
 \sum_{i,j=1}^{15}a_{ij}x_j\frac{\partial C}{\partial x_i}
     &=\lambda_3C+\left(\sum_{j=1}^{15}b_jx_j\right)Q
                                                        \label{eq:lie-cubic}.
\end{align}
These equations are linear in the $242$ unknowns
\[
 (a_{11},a_{12},\ldots,a_{15,15},\lambda_2,\lambda_3,b_1,\ldots,b_{15}).
\]
Here and below the order on the $a_{ij}$ is row-major.

We give the exact matrix recipe for the calculation.  Expand $Q$ and $C$ by the
matching formula above after the indicated substitution for $t_{44}$.  They have
$72$ and $222$ nonzero monomials, respectively.  In the column for $a_{ij}$ put
the coefficients of
\[
 x_j\partial_iQ+x_j\partial_iC.
\]
In the next two columns put the coefficients of $-Q$ and $-C$, and in the last
fifteen columns the coefficients of $-x_jQ$.  Rows are indexed by all monomial
exponent vectors that occur, in lexicographic order; their different total
degrees keep the quadratic and cubic equations separate.  The resulting integer
matrix $M$ has size $780\times242$.

The program \emph{Level-scheme stabilizer}, available through the workbench,
constructs precisely this matrix and records original row and column indices of
a $241\times241$ minor.  Its determinant is $681$ modulo the prime $1009$.
Thus $\operatorname{rank}_{\mathbb C}M\geq241$.  This is an exact integer-matrix
certificate: a nonzero determinant after reduction modulo a prime is nonzero
before reduction.  On the other hand, homogeneity gives the nonzero solution
\[
 A=\operatorname{id}_H,\qquad\lambda_2=2,\qquad\lambda_3=3,
 \qquad b_1=\cdots=b_{15}=0.
\]
Therefore the kernel has dimension exactly one, and every possible $A$ is a
scalar matrix.  Scalars preserve $I$, so this identifies the Lie algebra rather
than merely bounding it.

Finally let an algebraic torus act linearly and preserve $I$.  Decompose $H$ into
its character spaces.  If two distinct characters occurred, an integral
cocharacter of the torus would distinguish them.  Differentiating that
one-parameter subgroup over $\mathbb C$ would give a nonscalar element of the
Lie algebra just computed.  This is impossible.  All characters are equal, so
the torus acts by scalars.
\end{proof}

The theorem concerns the displayed scheme ideal.  It does not identify the
radical of that ideal, classify nonlinear automorphisms of its Fano schemes, or
show that arbitrary level spaces specialize to eigenspaces.  The diagonal
row-and-column torus used by Chan and Ilten
\cite[Section~2.2, Proposition~2.3 and Remark~2.5]{CI} acts on the
homogeneous permanental Fano scheme.  It moves the all-ones base point and does
not preserve the translated level equations here.

\section{Two restrictions on a possible six-dimensional space at order four}

The following restriction records explicitly where a small-case argument still
needs a classification.  Let $B$ be the polarization of $S_2$ on $H$, so that
$B(t,u)=S_2(t+u)-S_2(t)-S_2(u)$.  It is nondegenerate.

\begin{proposition}\label{prop:normal-rank-four}
Suppose $L\subseteq H$ is a six-dimensional level space of $4\times4$ matrices.
For each row and each column there is a linear functional annihilating its
projection of $L$ and taking value one on the all-ones vector.  Let $M\subseteq H^*$
be the span of any eight such functionals, restricted to $H$.  Identify $H^*$ with
$H$ by $B$.  The Gram form on $M$ has rank at most three.
\end{proposition}

\begin{proof}
A row projection of $L$ cannot contain the all-ones vector: otherwise it contains
its negative, and $J+t$ has a zero row for some $t\in L$, contrary to its permanent
being $24$.  Linear separation from a subspace supplies a functional vanishing
on the projection and nonzero on the all-ones vector; rescale it to value one.
The column argument is identical.

Put $r=\dim M$ and let $\rho$ be the rank of its induced Gram form.  The subspace
$U=M^\perp$ has dimension $15-r$ and radical
$M\cap M^\perp$ of dimension $r-\rho$.  Hence its largest totally isotropic
subspace has dimension
\[
 (r-\rho)+\left\lfloor\frac{15-2r+\rho}{2}\right\rfloor
       =\left\lfloor\frac{15-\rho}{2}\right\rfloor.
\]
The inclusion $L\subseteq U$ and the equality $S_2|_L=0$ give
$6\leq\lfloor(15-\rho)/2\rfloor$, so $\rho\leq3$.
\end{proof}

For explicit equations, put $D=\frac13J_4-I_4$.  On the full matrix space the
inverse of the polarization matrix is $D\otimes D$.  Passing to $H$ subtracts
$1/144$ from the pairing of any two coefficient vectors whose sums are one.
Thus normalized row functionals with coefficient vectors $a_i,a_k$ have pairing
\[
 D_{ik}\left(\frac13-a_i\cdot a_k\right)-\frac1{144}.
\]
The other entries follow from the same tensor formula.  Vanishing of the
$4\times4$ minors of this eight-generator Gram matrix is therefore a concrete
necessary condition on a six-dimensional candidate.  It does not yet classify
such candidates.

There is a separate obstruction to proving the bound only by isolating the
standard five-dimensional hook.  Its Grassmannian tangent space has ambient
dimension $5(15-5)=50$.  Exact differentiation of the layer equations gives ranks
$15$ for $S_2$, $33$ for $S_2,S_3$, and again $33$ after $S_4$.  Thus the tangent
space to the full level-space scheme there has dimension $17$.  The quartic has
zero first derivative at a hook because it would require three hook entries on
distinct rows and columns.  These are exact rational linear-algebra calculations,
with a reproducible coefficient construction in the workbench.  Higher-order
obstructions, rather than first-order isolation, are needed at this point.

